% 基本とドライバ関連
\usepackage{luatexja}
\usepackage{graphicx, xcolor}
\usepackage[top=25truemm,bottom=25truemm,left=20truemm,right=20truemm]{geometry}
\usepackage{wrapfig}
\usepackage{pdfpages}

\makeatletter
\def\caption@documentclass{elsarticle}% beamer, thesis
\newcommand{\tabcaption}[1]{\def\@captype{table}\caption{#1}}
\makeatother


% 数式(演算子など)のスペースを詰める
% =,→ 間の余白
% \thickmuskip=1.2\thickmuskip
% +,- 間の余白
% \medmuskip=0.8\medmuskip
% … などの装飾記号の余白
% \thinmuskip=0.8\thinmuskip
% 行列を詰める
% \arraycolsep=0.3\arraycolsep
% 数式の上下のスペースの変更
\AtBeginDocument{
	% \renewcommand{\abovedisplayskip}{3pt}
	% \renewcommand{\belowdisplayskip}{3pt}
	\abovedisplayskip=-4pt
	% \setlength{\abovedisplayshortskip}{0.3\abovedisplayshortskip}
	% \setlength{\belowdisplayskip}{0.3\belowdisplayskip}
	% \setlength{\belowdisplayshortskip}{0.3\belowdisplayshortskip}
}


% 参考資料系
\usepackage{here}
\usepackage{booktabs}
\usepackage{multirow}
\usepackage{multicol}
\usepackage[font=small, labelsep=quad]{caption}
\makeatletter
\let\MYcaption\@makecaption
\makeatother

\usepackage{subcaption}
\captionsetup{compatibility=false}	  % 必要に応じて

\makeatletter
\let\@makecaption\MYcaption
\makeatother

\usepackage[super]{cite}
\renewcommand\citeform[1]{[#1]}

% itembox用
\usepackage{ascmac}
\usepackage{fancybox}
% enumerate用
\renewcommand{\labelenumi}{(\arabic{enumi})}
\renewcommand{\labelenumii}{\alph{enumii})}
\renewcommand{\labelenumiii}{\roman{enumiii})}
% \renewcommand{\labelitemi}{*}


% フォント関連
\usepackage[no-math]{fontspec}
\usepackage[deluxe,expert,bold]{luatexja-preset} % hiragino-pron
% \usepackage{luatexja-fontspec}

% 数式
\usepackage{empheq}
\usepackage{amsmath,amsfonts,amssymb}
\usepackage{bm,ulem} % 太文字，下線
\usepackage{mathtools}
% \usepackage[warnings-off={mathtools-colon}]{unicode-math}
\numberwithin{equation}{section}
\mathtoolsset{showonlyrefs=true}

% bib関連
% \usepackage[backend=biber, style=numeric]{biblatex}
% \usepackage[backend=biber,url=false,style=ieee,sorting=nyt,doi=false]{biblatex}

% オリジナルコマンド
\renewcommand{\sp}{\vspace{2mm}}
\newcommand{\hsp}{\hspace{2mm}}
\newcommand{\f}[2]{\frac{#1}{#2}}
\newcommand{\wavecom}[2]{\underset{#1}{\uwave{#2}}}
\newcommand{\undlcom}[2]{\underset{#1}{\dashuline{#2}}}
\renewcommand{\(}{\left(}
\renewcommand{\)}{\right)}
\renewcommand{\[}{\left[}
\renewcommand{\]}{\right]}
\newcommand{\ol}[1]{\overline{#1}}
\newcommand{\sbsbsbsec}[1]{\vspace{2mm}\noindent\underline{#1}}

\makeatletter
\newcommand{\Kintou}[2]{\@tempcnta=0
	\@tfor\STR:=#2\do{\advance\@tempcnta\@ne}%
	\@tempcntb=\@tempcnta
	\leavevmode
	\hbox to #1\zw{\@tempcnta=0
		\@tfor\STR:=#2\do{\advance\@tempcnta\@ne
			\STR
			\ifnum\@tempcnta<\@tempcntb\hfill\fi}}}
\makeatother

%回路図
\usepackage{siunitx}
\usepackage{float}
\usepackage{tikz}
\usepackage{circuitikz}
\ctikzset{logic ports=ieee}
\usetikzlibrary{positioning}

% カルノー図 tikz使用
\usepackage{askmaps}
\usepackage{tikz -timing}
\usetikzlibrary{karnaugh}
\definecolor{LogisimKMapColor0}{HTML}{696969}
\definecolor{LogisimKMapColor1}{HTML}{4682B4}
\definecolor{LogisimKMapColor2}{HTML}{1E90FF}
\definecolor{LogisimKMapColor3}{HTML}{20B2AA}
\definecolor{LogisimKMapColor4}{HTML}{008B8B}
\definecolor{LogisimKMapColor5}{HTML}{006400}
\definecolor{LogisimKMapColor6}{HTML}{32CD32}
\definecolor{LogisimKMapColor7}{HTML}{FFD700}
\definecolor{LogisimKMapColor8}{HTML}{FF8C00}
\definecolor{LogisimKMapColor9}{HTML}{B22222}
\definecolor{LogisimKMapColor10}{HTML}{FF4500}
% \definecolor{LogisimKMapColor11}{HTML}{}
% \definecolor{LogisimKMapColor12}{HTML}{}
% \definecolor{LogisimKMapColor13}{HTML}{}
% \definecolor{LogisimKMapColor14}{HTML}{}
% \definecolor{LogisimKMapColor15}{HTML}{}

% ハイパーリンク関係
\usepackage{bookmark,xurl}
\urlstyle{same}
\hypersetup{
	unicode,
	bookmarksnumbered=true,
	hidelinks,
	final,
}

% ソースコード関連
\usepackage{listings, chngcntr}
\renewcommand{\lstlistingname}{コード}
\definecolor{Green}{HTML}{1D8991}
\definecolor{Red}{HTML}{DA103F}
\definecolor{Orange}{HTML}{F76C2B}
\definecolor{Purple}{HTML}{9847BF}
\definecolor{Gray}{HTML}{989290}
\definecolor{Pink}{HTML}{FDF0EC}
\definecolor{Black}{HTML}{333333}
% \lstdefinestyle{Default}{
\lstset{
	language={C},
	% morecomment=[s][\color{Orange}]{\ <}{\ },
	breaklines=true, % 行の折り返し
	breakindent=10pt, % 自動改行後のインデント量(デフォルトでは20[pt])
	showstringspaces=false, % 文字列の空白を表示
	numbers=left, % 行番号の指定
	numbersep=2pt,  % 行番号とプログラムの間隔
	stepnumber=1, % 行番号の間隔
	xrightmargin=0pt, % 右余白
	xleftmargin=20pt, % 左余白
	belowskip=7pt, % 下余白
	numberstyle={\scriptsize}, % 行番号の書体
	lineskip=-0.5ex,  % 行間の指定
	tabsize = 4, % タブの大きさ
	% 枠 "t"は上に線を記載, "T"は上に二重線を記載
	% 他オプション: leftline, topline, bottomline, lines, single, shadowbox
	frame=single,
	framexleftmargin=5mm,
	framexrightmargin=-2mm,
	frameround=tttt,
	% frame=single, % default=tb
	% frameまでの間隔(行番号とプログラムの間)
	framesep={3pt},
	framerule=3pt, % 枠線の太さ
	captionpos=b, % キャプションの場所("tb"ならば上下両方に記載)
	columns=[l]{fullflexible}, % 表示幅の指定
	mathescape=true, % 数式をエスケープ
	escapeinside={\@}{\@},
	%% 色，書体 %%
	% backgroundcolor={\color{Pink}}, % 背景色
	basicstyle={\ttfamily\small\color{Black}}, % プログラムの書体
	% identifierstyle={\color{Green}}, % 識別子の書体
	% identifierstyle={\bfseries},
	commentstyle={\itshape\color{Gray}}, % コメントの書体
	% keywordstyle={\color{Purple}}, % キーワードの書体
	% ndkeywordstyle={\color{Purple}}, % 予約語の書体
	% stringstyle={\color{Orange}}, % 文字列の書体
	% stringstyle={\bfseries},
}
% \lstnewenvironment{CCode}[1][]{
%   \lstset{style=Default, language={C},
%	 morecomment=[l][\color{Red}]{\#}, #1}
% }{}

\usepackage{tcolorbox}
\newtcbox{\codeinline}[1][]{
	% colback=Pink,
	% colframe=black,
	% colframe=Pink!0!Pink,
	boxrule=0pt,
	left=-2mm,right=0mm,top=0mm,bottom=0mm,
	box align=base,
	nobeforeafter,
	fontupper=\ttfamily
}
\newcommand{\code}[2]{\codeinline{\lstinline[#1]!#2!}}

% コンパイル時に図をとばす
\newif\iffigure
\figurefalse
\figuretrue

%数式に節番号を追加. jbookを使ったときはsection -> chapterへ変更.
\makeatletter
% 各種番号を"<節番号>.<番号>" へ
\renewcommand{\thefigure}{\thesection.\arabic{figure}}
\renewcommand{\thetable}{\thesection.\arabic{table}}
\renewcommand{\theequation}{\thesection.\arabic{equation}}

% 節が進むごとに図番号をリセットする
\@addtoreset{figure}{section}
\@addtoreset{table}{section}
\@addtoreset{equation}{section}
\makeatother

\AtBeginDocument{
	\counterwithin{lstlisting}{section}
}

\allowdisplaybreaks